\section{\texorpdfstring{Totalization,\allowbreak{} canonical triangles,\allowbreak{} and the vanishing-composition bound}{ canonical  and the vanishing-composition bound}}\label{proofs-3487-4161-totalization-canonical-triangles-and-the-vanishing-composition-bound}

Authority:\allowbreak{} the Stacks Project authors\textquotesingle{} \texttt{derived.\allowbreak{}tex},\allowbreak{} commit \texttt{a04446e5\allowbreak{}7ec1fbc2\allowbreak{}52a871af\allowbreak{}cec7752f\allowbreak{}b2807b14},\allowbreak{} bound in \texttt{INPUT\_\allowbreak{}BINDING.\allowbreak{}json}. Line numbers without a chapter prefix refer to that original file. Other chapter passages are the current public sources bound in the accompanying reading record. This is an editorial proof record. Source statements and their signs remain identifiable; no novelty is claimed.

\subsection{\texorpdfstring{1. Exact totalization,\allowbreak{} preserving both component conventions}{1. Exact  preserving both component conventions}}\label{proofs-3487-4161-exact-totalization-preserving-both-component-conventions}

The source uses commuting double-complex differentials and the full total differential \[
T(A)^n=\bigoplus_{p+q=n}A^{p,q},\qquad
D|_{A^{p,q}}=d_1^{p,q}+(-1)^p d_2^{p,q}.
\] The two diagonal squares of \(D^2\) vanish because \(d_1^2=d_2^2=0\). Its two mixed terms are \((-1)^{p+1}d_2^{p+1,q}d_1^{p,q}\) and \((-1)^pd_1^{p,q+1}d_2^{p,q}\); they cancel by the commuting-square condition. Equality on every summand inclusion proves the equality on the coproduct; no rearrangement of an infinite sum of morphisms is required.

For the first convention,\allowbreak{} with outer degree \(q\),\allowbreak{} the source shift is \([0,1]\). The exact functor comparison \(\xi_A:T(A[0,1])\to T(A)[1]\) multiplies the summand \(A^{p,q+1}\) by \((-1)^p\). This is the inverse of the original \(\gamma\) in Homology,\allowbreak{} and its inverse has the same component sign. On a horizontal differential,\allowbreak{} the source has coefficient \(1\) and the target has coefficient \(-1\),\allowbreak{} balanced by \((-1)^{p+1}=-(-1)^p\). On a vertical differential,\allowbreak{} both sides have coefficient \((-1)^{p+1}\) before applying that same sign. Thus this is an actual chain isomorphism,\allowbreak{} natural in all bigraded chain maps.

If \(h^{p,q}:A^{p,q}\to B^{p,q-1}\) is an outer homotopy,\allowbreak{} its total homotopy is \(H|_{A^{p,q}}=(-1)^ph^{p,q}\). The four terms of \(DH+HD\) on this summand are \[
(-1)^pd_1h+d_2h+(-1)^{p+1}hd_1+hd_2.
\] The first and third cancel because each \(h^q\) is a horizontal chain map. The other two are precisely \(f^{p,q}-g^{p,q}\). For a split sequence of these outer complexes,\allowbreak{} totalizing the actual sections and retractions gives a split sequence degree by degree. The connecting map on \(C^{p,q}\) is \[
\pi^{p+1,q}d_1s^{p,q}
+(-1)^p\pi^{p,q+1}d_2s^{p,q}
=(-1)^p\pi^{p,q+1}d_2s^{p,q}.
\] The omitted first term is zero for a reason:\allowbreak{} the section commutes with \(d_1\),\allowbreak{} and \(\pi s=0\). The remaining term is exactly \(\xi_A T(\delta)\). Hence images of the distinguished split triangles have the required boundary.

For the second convention,\allowbreak{} with outer degree \(p\),\allowbreak{} the shift is \([1,0]\). The comparison \(T(A[1,0])\to T(A)[1]\) is the identity on the reindexed summands. For an outer homotopy \(h^{p,q}:A^{p,q}\to B^{p-1,q}\),\allowbreak{} use \(H=h\). The four terms are \[
d_1h+(-1)^{p-1}d_2h+hd_1+(-1)^phd_2.
\] Here the two vertical terms cancel,\allowbreak{} and the horizontal pair gives \(f-g\). The connecting map of a totalized split sequence is \(\pi^{p+1,q}d_1s^{p,q}\):\allowbreak{} its vertical term vanishes because the sections are vertical chain maps. These are the two exactness claims at 3487–\allowbreak{}3532,\allowbreak{} with the full original comparisons retained. They agree with the previously published Homology review,\allowbreak{} rather than defining a replacement convention.

\subsection{2. DERIVED-NEW-013 and DERIVED-UNDER-004}\label{proofs-3487-4161-derived-new-013-and-derived-under-004}

For fixed \(X^\bullet\in\operatorname{Comp}(\mathcal A)\),\allowbreak{} the variable \(Y^\bullet\) belongs to \(\operatorname{Comp}(\mathcal B)\). Consequently the domain at 3572 is \(\operatorname{Comp}(\mathcal B)\). For fixed \(Y^\bullet\),\allowbreak{} the domain of the functor varying \(X^\bullet\),\allowbreak{} at 3574,\allowbreak{} is \(\operatorname{Comp}(\mathcal A)\). Those two source categories are interchanged in the printed proof. The two theorem assertions at 3555–\allowbreak{}3568 already have the correct domains. \textbf{DERIVED-NEW-013} repairs just the two domains in the proof.

Bilinearity supplies all needed identities:\allowbreak{} \(1_X\otimes h\) is the outer homotopy for the first functor,\allowbreak{} the horizontal chain-map identity follows from functoriality in the other variable,\allowbreak{} and the five biproduct splitting equations are preserved by \(X^p\otimes-\). The second variable-fixed construction uses \(h\otimes1_Y\) and the same equations. Thus both functors into the respective double-complex homotopy categories are exact with their literal outer shifts. Composing with §1 proves the totalized claims. The first totalized comparison has component \((-1)^p\) on \(X^p\otimes Y^{q+1}\); the second has component identity on \(X^{p+1}\otimes Y^q\). The direct receiving tensor lemmas in More on Algebra,\allowbreak{} Cohomology and Cohomology on Sites therefore keep their original conclusions and all flatness-related distinctions.

The same argument yields an editorial extension without requiring all countable coproducts. Let \(\mathcal E\) be an additive category. Inside either of the two homotopy categories of double complexes,\allowbreak{} take the full subcategory on those actual double complexes \(A\) for which every displayed diagonal coproduct \(\bigoplus_{p+q=n}A^{p,q}\) exists. This full subcategory is triangulated,\allowbreak{} and totalization on it is exact with the same comparisons.

Here is the closure and exactness proof. Diagonal sums of \(A[0,1]\) and \(A[1,0]\) are reindexings of the next diagonal sum of \(A\); the inverse shifts work in the same way. For a first-convention cone of \(f:A\to B\),\allowbreak{} the terms are \(B^{p,q}\oplus A^{p,q+1}\),\allowbreak{} so its diagonal sum exists and is canonically the finite biproduct \(T(B)^n\oplus T(A)^{n+1}\). For the second convention they are \(B^{p,q}\oplus A^{p+1,q}\),\allowbreak{} with the same conclusion. These identifications follow from the coproduct universal property,\allowbreak{} without discarding or identifying any labelled summands. The zero object and finite direct sums also remain in the subcategory. The already proved triangulated-subcategory criterion therefore applies.

In the first convention the isomorphism from the totalized outer cone to \(C(T(f))\) is identity on each \(B^{p,q}\) and \((-1)^p\) on each \(A^{p,q+1}\). The totalized outer-cone differential has upper-right entry \((-1)^pf\),\allowbreak{} source-part horizontal differential \(d_1\),\allowbreak{} and source-part vertical differential \(-(-1)^pd_2\). After the stated map,\allowbreak{} these become respectively \(f\),\allowbreak{} \(-d_1\),\allowbreak{} and \(-(-1)^pd_2\),\allowbreak{} the three entries required by the original cone of \(T(f)\). The changing sign between horizontal source indices accounts for \(-d_1\). The inclusion square is identity,\allowbreak{} and the last square is the exact comparison \(\xi_A\) on the source shift,\allowbreak{} with the retained negative projection. In the second convention the comparison is identity on both sets of summands:\allowbreak{} the outer cone already has source-part differential \(-d_1+(-1)^{p}d_2\) at new index \(p\),\allowbreak{} which is \(-d_1-(-1)^{p+1}d_2\) at its old source index. Its upper-right entry is \(f\). Thus totalization preserves distinguished cone triangles and is exact in both cases.

For a fixed tensor factor,\allowbreak{} restrict the variable-complex homotopy category to those objects for which the tensor double complex has these diagonal coproducts. Bilinearity and the same cone and shift calculation prove closure and exactness on that full triangulated subcategory. In particular,\allowbreak{} when both factors are bounded on the same side,\allowbreak{} every diagonal has finite support and finite biproducts suffice. For lower bounds \(p\geq a,q\geq b\),\allowbreak{} a fixed diagonal satisfies \(a\leq p\leq n-b\); for upper bounds it satisfies \(n-b\leq p\leq a\). These are the actual finite index intervals.

This is \textbf{DERIVED-UNDER-004},\allowbreak{} propagating the diagonal-coproduct calculation in \textbf{HOMOLOGY-UNDER-009} to exact functors and their domains. Existence of the sums remains essential. The Homology counterexample in finite-dimensional vector spaces,\allowbreak{} with identity horizontal pairs in bidegrees \((2i,-2i),(2i+1,-2i)\),\allowbreak{} has no degree-zero diagonal coproduct. It is excluded by the precise domain just defined. Also that domain need not be closed under arbitrary outer homotopy equivalences:\allowbreak{} the same example is outer contractible. No claim of strict fullness is made. No source theorem is silently replaced by this editorial extension.

\subsection{\texorpdfstring{3. Cohomology,\allowbreak{} localization and bounded models}{3.  localization and bounded models}}\label{proofs-3487-4161-cohomology-localization-and-bounded-models}

For a termwise split short exact sequence,\allowbreak{} its connecting morphism is \(\delta^n=\pi^{n+1}d_B^ns^n\). On a cycle,\allowbreak{} its image in \(H^{n+1}(A)\) is the snake-lemma boundary obtained by lifting along \(s^n\),\allowbreak{} applying \(d_B^n\),\allowbreak{} and retracting along \(\pi^{n+1}\). Independence follows because a different lift adds an \(A^n\)-component and therefore changes the result by an \(A\)-boundary. The general abelian-category assertion is the same pullback/\allowbreak{}kernel calculation in the referenced Homology lemma,\allowbreak{} not an assumption that a morphism from every test object can be lifted. Consequently \(H^0\) is homological and the source\textquotesingle s long sequence has boundary \(H^i(\delta)\) with its displayed sign.

\textbf{OCC-01499 /\allowbreak{} P02-E0027} and French \textbf{DERIVED-016} correctly request deletion of the unfinished reference placeholder at 3651; \textbf{MC-STK-ERR-0827} already does this,\allowbreak{} retaining both actual Homology references. French \textbf{DERIVED-017} matches \textbf{MC-STK-ERR-0828},\allowbreak{} the insertion of ``by'' at 3702. The functor being named is the unique factorization of the original \(H^0\).

The shifted kernel of \(H^0\) is exactly the acyclic complexes. For a distinguished triangle,\allowbreak{} the long sequence proves closure under cones and both shifts,\allowbreak{} and additivity proves closure under summands. The cone of a map is acyclic precisely when its cohomology maps are isomorphisms. Thus the localization denominators are exactly the quasi-isomorphisms,\allowbreak{} its kernel is exactly the acyclic subcategory,\allowbreak{} and its canonical shift comparison is identity. The quotient and bounded-subcategory constructions established earlier apply with these exact objects.

The size warning in the next remark is substantive. In the cited example,\allowbreak{} kernels and cokernels of an equivariant abelian-group map inherit each ordinal-indexed operator,\allowbreak{} and extending the shorter operator list by zeros gives the necessary common index. Finite biproducts do the same. This verifies the abelian structure used there. For the extensions of \(Z\) by \(Z\),\allowbreak{} an extension isomorphism has underlying matrix \(\begin{psmallmatrix}1&t\\0&1\end{psmallmatrix}\). If the nonzero operator of one extension occurs at ordinal \(\alpha\) and that of another at \(\beta\ne\alpha\),\allowbreak{} intertwining the operator at \(\alpha\) would require the product of that invertible matrix with \(\begin{psmallmatrix}0&1\\0&0\end{psmallmatrix}\) to be zero. It is not zero. Thus the distinct ordinals yield distinct extension classes. The use of their identification with derived Hom is the cited later Ext lemma,\allowbreak{} whose complete source-order review remains later; no new set-sized Hom claim is made for this large category.

For a Grothendieck abelian category the source invokes the specific K-injective existence theorem in Injectives. Its statement and the start of its construction were read here; its entire transfinite proof was not reaudited. The receiving comparison is exact:\allowbreak{} if \(I\) is K-injective and \(s:M\to N\) is a quasi-isomorphism,\allowbreak{} its acyclic cone and all shifts have zero Hom into \(I\). The long Hom sequence makes \(\operatorname{Hom}_K(N,I)\to\operatorname{Hom}_K(M,I)\) bijective. Every roof \(N\leftarrow M\to I\) therefore has a unique representative \(N\to I\); common refinements prove that this identification is independent of the roof. Composing with the selected quasi-isomorphism \(L\to I\) gives exactly the Hom comparison at 3746–\allowbreak{}3747,\allowbreak{} and its right side is a set.

The bounded models are obtained with their actual boundary objects:\allowbreak{} \[
(\tau_{\leq b}K)^i=
\begin{cases}K^i&i<b,\\\ker d_K^b&i=b,\\0&i>b,\end{cases}
\quad
(\tau_{\geq a}K)^i=
\begin{cases}0&i<a,\\\operatorname{coker}d_K^{a-1}&i=a,\\K^i&i>a.
\end{cases}
\] Their differentials are the original ones or their induced kernel and cokernel maps. Choose \(a<b\) outside the nonzero cohomology range. The two truncations commute on each degree,\allowbreak{} including the distinct boundary degrees,\allowbreak{} so their composite is the original \(N\) in the four-corner diagram. All four maps induce the identity on the surviving cohomology and zero-to-zero outside it.

For bounded-below complexes a left fraction \(s^{-1}f\) can have its common target replaced by its lower truncation,\allowbreak{} since that target is quasi-isomorphic to the bounded-below denominator source. This gives fullness. If the fraction is zero,\allowbreak{} a further denominator kills \(f\); lower-truncate its target again to give the same equality inside the bounded category. Additivity reduces equality of arbitrary fractions to this zero test,\allowbreak{} proving faithfulness. For bounded-above complexes use right fractions and upper truncation of their common source. For bounded complexes perform the bounded-below argument inside the bounded-above category; lower truncation preserves that upper bound. This supplies all three equivalences,\allowbreak{} with the kernel and essential surjectivity already verified.

\subsection{4. The canonical delta-functor and truncation triangles}\label{proofs-3487-4161-the-canonical-delta-functor-and-truncation-triangles}

For the original short exact sequence with maps \(a,b\),\allowbreak{} put \(q=(b,0):C(a)\to C\). Its kernel is identified with \(C(1_A)\) by \((x,y)\mapsto(a(x),y)\). The differential under that identification is \(\begin{psmallmatrix}d_A&1\\0&-d_A\end{psmallmatrix}\). The homotopy \(H^n(x,y)=(0,x)\) has \(dH+Hd=(x,y)\),\allowbreak{} so the kernel is contractible. The exact cohomology sequence makes \(q\) a quasi-isomorphism.

For a commuting map of short exact sequences \((f,g,h)\),\allowbreak{} the canonical cone map in degree \(n\) is \(c^n=\begin{psmallmatrix}g^n&0\\0&f^{n+1}\end{psmallmatrix}\). It satisfies \(q'c=hq\),\allowbreak{} \(p'c=f[1]p\). Hence,\allowbreak{} in the derived category,\allowbreak{} \[
(-p'(q')^{-1})h=-p'c q^{-1}=f[1](-pq^{-1}).
\] This proves the actual boundary naturality,\allowbreak{} including the shift. The signed cone triangle and the isomorphism \(q\) prove the distinguishedness. This is a further checked receiving use of DERIVED-NEW-007. All bounded variants preserve the cone\textquotesingle s one shift and finite sum. French \textbf{DERIVED-018} matches \textbf{MC-STK-ERR-0829},\allowbreak{} adding ``above'' at 3968 without changing this construction.

When all three vertical maps are quasi-isomorphisms,\allowbreak{} this same naturality gives the claimed triangle isomorphism. In the split case,\allowbreak{} the earlier homotopy inverse \(S=(s,-\delta)\) satisfies \(qS=1\),\allowbreak{} \(Sq\simeq1\),\allowbreak{} and \(-pS=\delta\). Thus the canonical boundary is exactly the image of the original split boundary. \textbf{OCC-01500 /\allowbreak{} P02-E0028},\allowbreak{} French \textbf{DERIVED-019},\allowbreak{} and French \textbf{DERIVED-020} match \textbf{MC-STK-ERR-0830},\allowbreak{} \textbf{0831},\allowbreak{} the two singular-sequence copyedits.

For the first truncation triangle,\allowbreak{} write \(U=K/\tau_{\leq a}K\). The natural map \(U\to\tau_{\geq a+1}K\) has kernel concentrated in degrees \(a,a+1\),\allowbreak{} with both terms \(\operatorname{im}d_K^a\) and differential identity under the canonical image/\allowbreak{}coimage identification. Hence the kernel is acyclic and the map is a quasi-isomorphism. The canonical delta-functor triangle transports across that precise map. Applying the same construction to \(\tau_{\leq a+1}K\) leaves only the quotient cohomology \(H^{a+1}(K)[-a-1]\). For the lower truncations,\allowbreak{} the kernel of \(\tau_{\geq a}K\to\tau_{\geq a+1}K\) has term \(K^a/\operatorname{im}d^{a-1}\) at \(a\),\allowbreak{} term \(\operatorname{im}d^a\) at \(a+1\),\allowbreak{} and the induced surjection between them. Inclusion of its kernel \(H^a(K)[-a]\) is a quasi-isomorphism. This proves the third triangle and its final shift \([-a+1]\),\allowbreak{} with all original quotient contributions retained.

\subsection{\texorpdfstring{5. DERIVED-UNDER-005:\allowbreak{} an explicit bound and derived morphisms}{5.  an explicit bound and derived morphisms}}\label{proofs-3487-4161-derived-under-005-an-explicit-bound-and-derived-morphisms}

The induction at 4154–\allowbreak{}4160 applies its one-step case after shifting the cohomological cutoff. Its full form is the following. Let \(f_j:K_j\to K_{j+1}\),\allowbreak{} \(0\leq j<n\),\allowbreak{} be morphisms in \(D(\mathcal A)\),\allowbreak{} and fix an integer \(b\). Assume \[
H^i(K_0)=0\quad(i>b),\qquad H^{b-j}(f_j)=0\quad(0\leq j<n).
\] Then the composite factors through \(\tau_{\leq b-n}K_n\to K_n\). The intermediate complexes need not have any boundedness assumption. This includes arbitrary derived morphisms,\allowbreak{} not only the chain maps in the source statement.

We prove the needed one-step assertion without assuming a representative exists on the original pair of complexes. If \(X\) has cohomology zero above \(m\),\allowbreak{} represent \(f:X\to Y\) by a roof \(X\xleftarrow{s}Z\xrightarrow{g}Y\). The map \(\tau_{\leq m}Z\to Z\) is a quasi-isomorphism. The restricted chain map factors through \(\tau_{\leq m}Y\),\allowbreak{} because its degree-\(m\) component carries cycles to cycles. The composite \(\tau_{\leq m}Z\to H^m(Y)[-m]\) equals the canonical projection to \(H^m(Z)[-m]\) followed by \(H^m(g)\). If \(H^m(f)=0\),\allowbreak{} then \(H^m(g)=H^m(f)H^m(s)=0\). Thus that composite is zero as a chain map. Applying the exact representable Hom sequence to the canonical triangle \[
\tau_{\leq m-1}Y\to\tau_{\leq m}Y\to H^m(Y)[-m]
\] gives the required lift \(X\to\tau_{\leq m-1}Y\).

These truncation constructions are natural on the derived category. On complexes their kernel/\allowbreak{}cokernel definitions are functorial and preserve quasi-isomorphisms by the displayed cohomology formulas. They also preserve homotopies:\allowbreak{} at the upper cutoff restrict the last homotopy component to the kernel; at the lower cutoff project its first nonzero target component to the cokernel. The boundary equations are exactly the original \(f-g=dh+hd\),\allowbreak{} restricted or projected. The universal property of localization therefore supplies the induced functors and their canonical maps used above.

Starting with \(u_0:K_0\to\tau_{\leq b}K_0\),\allowbreak{} the inverse of the canonical quasi-isomorphism,\allowbreak{} suppose the first \(j\) maps factor as \(u_j:K_0\to\tau_{\leq b-j}K_j\) followed by the canonical inclusion. Apply the one-step assertion at \(m=b-j\) to \[
\tau_{\leq b-j}K_j\longrightarrow K_j\xrightarrow{f_j}K_{j+1}.
\] Its degree-\(b-j\) cohomology map vanishes by hypothesis. It therefore factors through \(\tau_{\leq b-j-1}K_{j+1}\). Composing that factorization with \(u_j\) constructs \(u_{j+1}\). Induction proves the stated result. For the source\textquotesingle s \(n=2\) step this is the application at cutoff \(-1\),\allowbreak{} equivalently the cutoff-zero case on complexes shifted by \([-1]\). It is an implicit shift in the source,\allowbreak{} not a failure of its conclusion.

The dual assertion says that for arrows \(K_n\to\cdots\to K_0\),\allowbreak{} with \(H^i(K_0)=0\) below \(a\) and \(H^{a+j}(K_{j+1}\to K_j)=0\),\allowbreak{} the composite factors through \(K_n\to\tau_{\geq a+n}K_n\). To obtain it,\allowbreak{} reverse arrows and degrees:\allowbreak{} a cochain complex in \(\mathcal A^{op}\) associated to \(K\) has degree \(i\) object \(K^{-i}\) and differential the opposite of the original map \(d_K^{-i-1}\). Cohomology in degree \(i\) is the opposite of \(H^{-i}(K)\); quasi-isomorphisms and their roofs are reversed,\allowbreak{} and upper truncation becomes lower truncation under this operation. Applying the just-proved factorization with bound \(-a\) and reversing it back gives exactly the claimed arrow and indices.

Now suppose additionally that \(H^i(K_n)=0\) for \(i<a\) and \(n>b-a\). The receiving truncation \(\tau_{\leq b-n}K_n\) is acyclic,\allowbreak{} hence zero in the derived category. The composite is zero. In particular,\allowbreak{} for a complex with cohomology in \([a,b]\),\allowbreak{} a composite of \(b-a+1\) endomorphisms that each act by zero on all cohomology is zero. These endomorphisms form an ideal,\allowbreak{} since cohomology is an additive functor. Its \((b-a+1)\)-st power is therefore zero. If \(u\) is one of them,\allowbreak{} the complete finite formula \[
(1-u)^{-1}=1+u+u^2+\cdots+u^{b-a}
\] follows by multiplying on either side and retaining the terminal term \(-u^{b-a+1}=0\). No global dimension or existence of projective/\allowbreak{}injective resolutions is used for this consequence.

This is \textbf{DERIVED-UNDER-005}. It is a proved corollary and extension of the original factorization argument,\allowbreak{} not a claim of a new theorem in the literature. Its source theorem stays unchanged.

\subsection{6. Propagation and DERIVED-NEW-014}\label{proofs-3487-4161-propagation-and-derived-new-014}

In \texttt{more-\allowbreak{}algebra.\allowbreak{}tex},\allowbreak{} current 23390–\allowbreak{}23402 (original 22711–\allowbreak{}22723)\allowbreak{},\allowbreak{} the earlier \textbf{MC-STK-ERR-0959} corrected both reversed composites to \(\psi\varphi\). The reports \textbf{OCC-01806 /\allowbreak{} P02-E0339},\allowbreak{} \textbf{OCC-02517 /\allowbreak{} P02-E1050},\allowbreak{} and \textbf{OCC-12698 /\allowbreak{} MORE-ALGEBRA-G-016} concern that order and do not ask to correct the receiving truncation. Their original report texts were read here; this is a propagation check,\allowbreak{} not a new whole-chapter report closure.

With the original labels \(K_1\xrightarrow\varphi K_2
\xrightarrow\psi K_3\),\allowbreak{} §5 applies at \(b=0,n=2\). It factors the composite through \textbf{\(\tau_{\leq-2}K_3\)}. The proof currently writes \(K_2\). Since both truncations vanish,\allowbreak{} the displayed existence of some factorization through zero is not a counterexample to the theorem. Nevertheless the cited argument\textquotesingle s receiving object is \(K_3\); \textbf{DERIVED-NEW-014} corrects that index in the proof and retains the earlier order corrections. The theorem itself needs only the source upper bound on \(K_1\),\allowbreak{} the final lower bound on \(K_3\),\allowbreak{} and the two specified vanishing cohomology maps; §5 proves this stronger editorial statement without requiring \(K_2\) to have amplitude \([-1,0]\).

The other direct uses were checked as receiving applications:\allowbreak{}

\begin{itemize}
\tightlist
\item
  In \texttt{cotangent.\allowbreak{}tex:\allowbreak{}2250–\allowbreak{}2288},\allowbreak{} the original \(c=k^2\) is retained. There are \(k\) transitions,\allowbreak{} each lowering the ideal index by \(k\),\allowbreak{} and each kills degrees \(0,-1,\ldots,-k+1\). The composite factors through the target cutoff \(-k\),\allowbreak{} exactly as asserted. For \(k=0\) the assertion is only the existing vanishing in positive cohomological degrees; no zero-fold composition is claimed to vanish.
\item
  In \texttt{local-\allowbreak{}cohomology.\allowbreak{}tex:\allowbreak{}3464–\allowbreak{}3484},\allowbreak{} the truncated dualizing complex has amplitude \([-d+1,0]\). Multiplication by the chosen \(f\) is zero on its cohomology,\allowbreak{} so its \(d\)-th power is zero. For \(d=0\) that truncated object is already zero,\allowbreak{} consistent with the conclusion.
\item
  In \texttt{local-\allowbreak{}cohomology.\allowbreak{}tex:\allowbreak{}4375–\allowbreak{}4393},\allowbreak{} the homological cone of the augmentation of the Koszul complex occupies degrees \(0\) through \(t\). Its cochain amplitude is therefore contained in \([-t,0]\). If the recorded \(a_1^c\) kills its cohomology,\allowbreak{} the explicit sufficient exponent is \(c(t+1)\). This proves the receiving step that enlarges \(c\); it does not claim to optimize the final ideal exponent in the whole lemma.
\item
  In \texttt{local-\allowbreak{}cohomology.\allowbreak{}tex:\allowbreak{}4685–\allowbreak{}4701},\allowbreak{} the \(t\) transitions kill negative degrees \(-t,\ldots,-1\). Apply the dual assertion with \(a=-t\) to the truncated objects to factor through the degree-zero quotient \(M/I^{n+tc'}M\) of the source. The displayed transitions at index \(n+tc'\) and the later index \(n-c\) are retained; no change of the constants or source objects is made.
\end{itemize}

These are bounded checks of the use of the factorization lemma. They do not certify every surrounding argument or report in those chapters. The broader joint synthesis remains after the core review queue.
